\chapter{Currents and compactness}\label{sec:currents}
Currents retain orientation and multiplicity. Their weak topology is defined by testing against differential forms. We use \cite{Sim18}, Chapter 6, and compare the overview in \cite{Fle16}, Sections 4--5.

\begin{definition}[Current and boundary]\label{def:current}
Let $\mathcal{D}^m(U)=C_c^\infty(U;\Lambda^m(\mathbb{R}^n)^*)$ be the space of smooth compactly supported $m$-forms, with its usual test-function topology. An \emph{$m$-current} is a continuous linear functional $T:\mathcal{D}^m(U)\to\mathbb{R}$. For $m\geq1$ its boundary is the $(m-1)$-current
\[
 \partial T(\eta)=T(d\eta),\qquad\eta\in\mathcal{D}^{m-1}(U).
\]
For $m=0$ set $\partial T=0$. Weak convergence $T_j\rightharpoonup T$ means $T_j(\omega)\to T(\omega)$ for every $\omega\in\mathcal{D}^m(U)$. See \cite{Sim18}, Chapter 6, Section 2, pp.~159--161.
\end{definition}

\begin{proposition}[Boundary squared and weak continuity]\label{prop:boundary-current}
Every current satisfies $\partial(\partial T)=0$. Moreover, $T_j\rightharpoonup T$ implies $\partial T_j\rightharpoonup\partial T$.
\end{proposition}
\begin{proof}
For an admissible test form $\eta$, $\partial^2T(\eta)=T(d^2\eta)=0$. For a test form $\zeta$ of degree $m-1$, its exterior derivative is again smooth and compactly supported, so
$\partial T_j(\zeta)=T_j(d\zeta)\to T(d\zeta)=\partial T(\zeta)$.
The assertions in dimension zero follow from the convention in Definition~\ref{def:current}.
\end{proof}

\begin{definition}[Mass and the comass convention]\label{def:mass-current}
The comass of an $m$-covector $\alpha$ is
\[
 \|\alpha\|_* =\sup\{|\langle\alpha,v_1\wedge\cdots\wedge v_m\rangle|:
             v_1,\ldots,v_m\text{ orthonormal}\}.
\]
For $m=0$ it is the absolute value. For an open $W\subset U$ define
\[
 \mathbf{M}_W(T)=\sup\{T(\omega):\omega\in\mathcal{D}^m(W),\
                 \sup_x\|\omega(x)\|_*\leq1\}.
\]
Write $\mathbf{M}(T)=\mathbf{M}_U(T)$. Test forms on $W$ are extended by zero to $U$.

We use the comass convention of \cite{Fed69}, Section 1.8, pp.~38--40. Simon initially uses the Euclidean norm on covectors in his mass definition, then discusses this distinction in \cite{Sim18}, Chapter 6, Remark 2.7, p.~160. Both conventions agree for the rectifiable currents in Definition~\ref{def:integral-current}; they need not agree for arbitrary currents. We will not identify their numerical values without this qualification.
\end{definition}

\begin{proposition}[Lower semicontinuity of mass]\label{prop:mass-lsc}
If $T_j\rightharpoonup T$ in $U$, then for every open $W\subset U$,
\[
 \mathbf{M}_W(T)\leq\liminf_j\mathbf{M}_W(T_j).
\]
\end{proposition}
\begin{proof}
For each test form $\omega$ in the supremum defining $\mathbf{M}_W(T)$,
$T(\omega)=\lim_jT_j(\omega)\leq\liminf_j\mathbf{M}_W(T_j)$.
Take the supremum over $\omega$. This proof applies even if the limiting lower bound is infinite. Compare \cite{Sim18}, Chapter 6, Section 2.
\end{proof}

\begin{definition}[Integer rectifiable and integral currents]\label{def:integral-current}
An $m$-current $T$ in $U$ is \emph{locally integer rectifiable} if
\[
 T(\omega)=\int_M\theta(x)\langle\omega(x),\tau(x)\rangle\,d\mathcal{H}^m(x),
\]
where $M\subset U$ is $\mathcal{H}^m$-measurable and countably $m$-rectifiable, $\theta:M\to\{1,2,\ldots\}$ is measurable with
$\int_{M\cap K}\theta\,d\mathcal{H}^m<\infty$ for every compact $K\subset U$, and $\tau(x)$ is a measurable unit simple $m$-vector orienting $T_xM$ almost everywhere. For $m=0$ this means a locally finite integer combination of point evaluations. It is \emph{locally integral} if both $T$ and $\partial T$ are locally integer rectifiable. It is integral with finite mass and boundary mass if also $\mathbf{M}(T)+\mathbf{M}(\partial T)<\infty$.

For the displayed representation,
$\mathbf{M}_W(T)=\int_{M\cap W}\theta\,d\mathcal{H}^m$.
The mass formula and representation are cited facts; see \cite{Sim18}, Chapter 6, Definition 3.1 and Remark 3.2(3), pp.~172--173. Local finiteness and integrability are taken in the ambient open set, not merely on compact subsets of $M$.
\end{definition}

\begin{theorem}[Federer--Fleming compactness]\label{thm:integral-compactness}
Let $1\leq m<n$. Suppose that $T_j$ are integral $m$-currents in $\mathbb{R}^n$, their supports lie in one compact set $K$, and
\[
 \sup_j\bigl(\mathbf{M}(T_j)+\mathbf{M}(\partial T_j)\bigr)<\infty.
\]
Then a subsequence converges weakly to an integral current $T$ supported in $K$, and
\[
 \mathbf{M}(T)\leq\liminf_k\mathbf{M}(T_{j_k}),\qquad
 \mathbf{M}(\partial T)\leq\liminf_k\mathbf{M}(\partial T_{j_k}).
\]
\emph{Cited compactness and closure theorem; proof to be added.} See \cite{Sim18}, Chapter 6, Theorem 3.15, p.~175, with its proof in Section 8, and \cite{Fle16}, Sections 4--5, pp.~5--7. The preservation of integer rectifiability is the substantial conclusion. The inequalities then follow from Propositions~\ref{prop:boundary-current} and~\ref{prop:mass-lsc}.
\end{theorem}

\chapter{Polyhedral and flat chains}\label{sec:flat-chains}
The flat norm makes boundary continuous while allowing thin fillings to become small. We begin with integer coefficients and polyhedral chains, following \cite{WhiCho14}, Sections 2.1--2.3. General coefficient groups will require separate hypotheses and are not silently included in the statements here.

\begin{definition}[Integer polyhedral chains]\label{def:polyhedral-chain}
The group $\mathcal{P}_m(\mathbb{R}^n;\mathbb{Z})$ consists of finite integer sums of oriented compact convex $m$-dimensional polyhedra, modulo subdivision and reversal of orientation. After a common subdivision, a chain can be written $P=\sum_i a_i[\sigma_i]$ with pairwise disjoint relative interiors and with repeated pieces combined. Its mass and boundary are
\[
 \mathbf{M}(P)=\sum_i|a_i|\mathcal{H}^m(\sigma_i),\qquad
 \partial P=\sum_i a_i\partial[\sigma_i],
\]
where faces carry their induced orientations. For $m=0$ use counting measure and set the boundary to zero. These constructions respect subdivision, and $\partial^2P=0$. See \cite{WhiCho14}, Definition 2.4 and equation (2.1), pp.~4--5.
\end{definition}

\begin{definition}[Flat norm and completion]\label{def:flat-norm}
For $P\in\mathcal{P}_m(\mathbb{R}^n;\mathbb{Z})$, define
\[
 \mathcal{F}(P)=\inf\{\mathbf{M}(R)+\mathbf{M}(S):
 R\in\mathcal{P}_m(\mathbb{R}^n;\mathbb{Z}),\
 S\in\mathcal{P}_{m+1}(\mathbb{R}^n;\mathbb{Z}),\ P=R+\partial S\}.
\]
When $m=n$, the group of $(n+1)$-chains is zero. This flat norm is symmetric, subadditive, and vanishes only at zero, so $d(P,Q)=\mathcal F(P-Q)$ is a translation-invariant metric; a proof is given in Proposition \ref{prop:polyhedral-flat-nondegeneracy}. For a compact set $K$, let $\mathcal{P}_m(K)$ be the chains supported in $K$, with the ambient flat norm. The compactly supported flat chains are the union over compact $K$ of the metric completions of $\mathcal{P}_m(K)$, using the natural inclusions. This is the convention of \cite{WhiCho14}, equation (2.2) and the discussion on pp.~5--6. The fillings in the displayed infimum are not required to lie in $K$.
\end{definition}

\begin{proposition}[Boundary is continuous in the flat norm]\label{prop:flat-boundary}
For every integer polyhedral chain $P$,
\[
 \mathcal{F}(P)\leq\mathbf{M}(P),\qquad
 \mathcal{F}(\partial P)\leq\mathcal{F}(P).
\]
The boundary therefore extends continuously to the completions in Definition~\ref{def:flat-norm}. See \cite{WhiCho14}, Section 2.3, pp.~6--8.
\end{proposition}
\begin{proof}
Take $R=P$ and $S=0$ for the first inequality. If $P=R+\partial S$, then $\partial P=\partial R$. In the definition of $\mathcal{F}(\partial P)$ use residual chain zero and filling $R$. Thus
$\mathcal{F}(\partial P)\leq\mathbf{M}(R)\leq\mathbf{M}(R)+\mathbf{M}(S)$.
Take the infimum over decompositions of $P$. Applying the same estimate to differences gives the continuous extension on each completion.
\end{proof}

\begin{example}[Thin boundaries and loss of mass]\label{ex:thin-rectangle}
In $\mathbb{R}^2$ let $Q_\varepsilon$ be the oriented rectangle $[0,1]\times[0,\varepsilon]$, with $\varepsilon>0$. Then
\[
 \mathcal{F}(\partial Q_\varepsilon)\leq\mathbf{M}(Q_\varepsilon)=\varepsilon,
 \qquad \mathbf{M}(\partial Q_\varepsilon)=2+2\varepsilon.
\]
Thus $\partial Q_\varepsilon$ converges to zero in the flat norm as $\varepsilon\downarrow0$, although its mass tends to $2$. The first estimate uses $Q_\varepsilon$ as a filling and the second sums its four edge lengths. This is the example in \cite{WhiCho14}, Example 2.7, p.~6. Flat convergence does not imply convergence of mass.
\end{example}
